\section{Introduction}
\label{sec:introduction}
\subsection{Background and motivation}
Image classification assigns a category to an image. For a product catalogue,
this could help organize new items into clothing and accessory groups.
Reliable classification depends on more than the model: the data split,
preprocessing and evaluation procedure also affect how a result should be
interpreted.

Fashion-MNIST provides a manageable setting for studying this process
\cite{fresh_fashion}. Its small grayscale images make it practical to train
several models, while visually similar categories create a useful recognition
challenge. A shirt and a pullover, for example, may have similar outlines even
though they belong to different labels.

Different architectures use image information in different ways. A Linear
classifier combines pixel values directly. An MLP learns nonlinear combinations
of pixels. A CNN learns local spatial patterns, an LSTM reads an ordered sequence
of rows, and a Transformer combines information across image patches.
Comparing these approaches helps explain what each design contributes in a
small, controlled experiment.

\subsection{Problem statement}
The task is supervised, single-label classification into ten categories.
A raw image is $X\in\{0,\ldots,255\}^{28\times28}$ and its label is
$y\in\{0,\ldots,9\}$. After preprocessing $T$, a model with parameters $\theta$
produces ten scores:
\[
f_\theta(T(X))=z\in\mathbb{R}^{10},
\qquad \hat y=\arg\max_{k\in\{0,\ldots,9\}}z_k.
\]
The prediction describes the complete image. The study does not locate objects
within an image or assign several labels to one image.

The main question is how the five implemented configurations compare when they
use the same training examples, normalization and epoch budget. We examine both
overall performance and the types of mistakes hidden by an aggregate score.

\subsection{Objectives and scope}
The project has four objectives:
\begin{enumerate}
\item Inspect the dataset and establish a reproducible preprocessing and split protocol.
\item Implement and compare Linear, MLP, CNN, LSTM and Transformer classifiers.
\item Evaluate accuracy, macro-F1, learning behavior and class-specific errors.
\item Relate predictive quality to parameter count and documented runtime measurements.
\end{enumerate}

The experiments use Fashion-MNIST only, with no pretrained weights, external
training data or augmentation. Each model is trained once using seed 42 and ten
epochs. The comparison concerns these particular configurations; architecture,
parameter count and computation are not independently controlled.

\subsection{Report approach}
Section~2 describes the data, exploratory findings and preprocessing.
Section~3 explains the models and the common training procedure.
Section~4 compares the saved results and runtime measurements.
Section~5 examines class-level and image-level errors.
Sections~6 and 7 discuss the limitations and conclusions.
The appendices provide the numerical class results, additional confusion matrices
and reproduction instructions.

The analysis uses the existing completed runs. Their saved predictions are used
to recompute accuracy, class scores and confusion counts. Validation selects
the checkpoint; the reported test results are kept separate from that selection.
